% Section 13: The Lie derivative and Cartan's formula} | Part: 2 | Depends on: see docs/REWRITE_PLAN.md | Status: todo | Original lines: 2146-2172
%!TEX root = ../../../main.tex
% =====================================================================
\section{The Lie derivative and Cartan's formula}\label{sec:lie}
% =====================================================================

How does a field change from point to point? For a function the question is easy; for vectors and forms we must compare objects living in different
vector spaces.
\begin{intuition}[{the problem}]
For a function $f$ it is easy to compare $f(p)$ and $f(q)$ at nearby points: they are both numbers. For a vector or tensor field, $v_p\in T_pM$ and $v_q\in T_qM$ live in \emph{different} vector spaces: to subtract them we need a rule to \textbf{identify $T_pM$ with $T_qM$}. There are two natural rules, giving two derivatives:
\begin{center}
\begin{tabular}{@{}p{3.6cm}p{5.6cm}p{5.2cm}@{}}
\toprule
 & identification of $T_pM\to T_{q}M$ & extra data needed\\ \midrule
\textbf{Lie derivative} $\Lie_X$ & differential of the flow of $X$ & none (but depends on $X$ \emph{near} $p$)\\
\textbf{Covariant derivative} $\nabla_X$ & parallel transport along the flow line & a \emph{connection} $\nabla$\\ \bottomrule
\end{tabular}
\end{center}
\end{intuition}


\begin{remark}[a third derivative]
A differential form admits a third derivative, the exterior derivative $\dd$ (Section~\ref{sec:dform}), which needs \emph{no identification of tangent spaces at all}: it
uses only the antisymmetrised partial derivatives, and is therefore independent of both a vector field (unlike $\Lie_X$) and a connection (unlike $\nabla$).
Cartan's formula below relates it to the Lie derivative, and Section~\ref{sec:synth3} compares all three.
\end{remark}
% Section 13.1: Vector fields and their flows | Part: 2 | Depends on: sec:fields, sec:interior | Status: todo | Original lines: 2173-2180
%!TEX root = ../../../main.tex
\subsection{Vector fields and their flows}
\begin{formal}[{flow of a vector field}]
A vector field $X\in\Xfield(M)$ is equivalent to an ODE for curves: $\dot\gamma(t)=X_{\gamma(t)}$, \ie in a chart $\dot x^i=X^i(x)$. By existence and uniqueness of ODE solutions $X$ defines a \emph{flow} $\phi_t$ (defined for small $t$ near each point): $\phi_t(x)$ is the solution at time $t$ with initial condition $x$. Properties: $\phi_0=\id$, $\phi_{t_1}\circ\phi_{t_2}=\phi_{t_1+t_2}$ (where defined), $\phi_t^{-1}=\phi_{-t}$, $\frac{\dd}{\dd t}\phi_t(x)=X_{\phi_t(x)}$, and to first order $\phi_\varepsilon(x)=x+\varepsilon X(x)+O(\varepsilon^2)$.
\end{formal}
\begin{intuition}[{wind blowing a leaf}]
$X$ is a stationary wind map; $\phi_t$ moves every point for time $t$ following the wind. The family $t\mapsto\phi_t$ is a one-parameter group of local diffeomorphisms: $X$ is its \emph{infinitesimal generator}.
\end{intuition}

% Section 13: Definition and properties} | Part: 2 | Depends on: see docs/REWRITE_PLAN.md | Status: todo | Original lines: 2181-2229
%!TEX root = ../../../main.tex
\subsection{Definition and properties}
\begin{formal}[{Lie derivative}]
For $\tau\in\TF{\ell}{k}(M)$ and $X\in\Xfield(M)$ with flow $\phi_\varepsilon$,
\[
(\Lie_X\tau)_p=\frac{\dd}{\dd\varepsilon}\Big|_{\varepsilon=0}(\phi_\varepsilon^*\tau)_p=\lim_{\varepsilon\to0}\frac{(\phi_\varepsilon^*\tau)_p-\tau_p}{\varepsilon}\ \in\TT{\ell}{k}(T_pM),
\]
\ie $(\phi_\varepsilon^*\tau)_p=\tau_p+\varepsilon(\Lie_X\tau)_p+O(\varepsilon^2)$. Properties:
\begin{itemize}[nosep]
\item $\Lie_Xf=X(f)$ for functions.
\item (Leibniz) $\Lie_X(\tau_1\otimes\tau_2)=\Lie_X\tau_1\otimes\tau_2+\tau_1\otimes\Lie_X\tau_2$; it ``differentiates every slot'': $\Lie_X(\tau(v_1,\dots))=(\Lie_X\tau)(v_1,\dots)+\tau(\Lie_Xv_1,\dots)+\dots$
\item For vector fields, $\Lie_XY=[X,Y]=X\circ Y-Y\circ X$ (the \emph{Lie bracket}), \ie $[X,Y]f=X(Yf)-Y(Xf)$.
\item $\Lie_X(\dd f)=\dd(\Lie_Xf)$.
\item In coordinates $X=X^j\partial_j$:\ $\Lie_X\partial_i=-\partial_iX^j\,\partial_j$,\quad $\Lie_X\dd x^i=\partial_jX^i\,\dd x^j$, and for a metric
\begin{equation}\label{eq:liemetric}
(\Lie_Xg)_{ab}=X^k\partial_kg_{ab}+g_{kb}\,\partial_aX^k+g_{ak}\,\partial_bX^k .
\end{equation}
\end{itemize}
\end{formal}

\begin{intuition}[{``how does $\tau$ change when I drag it with the flow, compared with the original?''}]
You push $\tau$ along the flow of $X$ for a tiny time (like a sheet of rubber being dragged by the water) and compare with what $\tau$ actually is there. $\Lie_X\tau=0$ means $\tau$ is \textbf{invariant} under the flow: it is a symmetry. The bracket $[X,Y]$ measures the failure of the flows of $X$ and $Y$ to commute: go $\varepsilon$ along $X$, then $\varepsilon$ along $Y$, then back along $X$ and $Y$: you end up displaced by $\varepsilon^2[X,Y]$ (a rectangle that does not close).
\end{intuition}

\begin{pitfall}[{$\Lie_X$ is \emph{not} a connection}]
$\Lie_{fX}Y=f\,\Lie_XY-Y(f)\,X\neq f\,\Lie_XY$. The Lie derivative needs the flow of $X$ near the point (the derivatives of $X$ appear in it, \eg $\partial_aX^k$ in \eqref{eq:liemetric}), not just the value $X_p$. A covariant derivative, in contrast, is $C^\infty(M)$-linear in $X$.
\end{pitfall}

\begin{example}[{a Lie bracket}]
On $\R^2$ take $X=\partial_x$, $Y=x\,\partial_y$. For any $f$: $[X,Y]f=\partial_x(x\,\partial_yf)-x\,\partial_y\partial_xf=\partial_yf$, so $[\partial_x,x\partial_y]=\partial_y\neq0$: moving first in $x$ and then in $y$ is not the same as in the opposite order because the $y$-speed depends on $x$.
\end{example}

\begin{formal}[{invariant tensors and Killing fields}]
$\tau$ is \emph{invariant under $X$} if $\Lie_X\tau=0$, equivalently $\phi_t^*\tau=\tau$ for all $t$. If a metric satisfies $\Lie_Xg=0$, then $X$ is called an \emph{infinitesimal isometry} or \textbf{Killing vector field}; its flow consists of (local) isometries.
\end{formal}

\begin{example}[{Killing field and conformal Killing field of the plane in polar coordinates}]
With $g=\mathrm{diag}(1,r^2)$ in polar coordinates and formula \eqref{eq:liemetric}:
\begin{itemize}[nosep]
\item $X=\partial_\varphi$ (rotations, $X^r=0$, $X^\varphi=1$): all $\partial_aX^k=0$ and $X^k\partial_kg_{ab}=\partial_\varphi g_{ab}=0$, hence $\Lie_Xg=0$: \textbf{Killing}. (Its flow is the rotation $\phi_\varepsilon(r,\varphi)=(r,\varphi+\varepsilon)$, which clearly preserves lengths.)
\item $X=r\,\partial_r$ (dilations $\phi_\varepsilon(x)=e^\varepsilon x$): $(\Lie_Xg)_{rr}=r\partial_r(1)+2g_{rr}\partial_rX^r=2$, $(\Lie_Xg)_{\varphi\varphi}=r\partial_r(r^2)=2r^2$, so $\Lie_Xg=2g$: not Killing but \emph{conformal} (lengths are rescaled by $e^{2\varepsilon}$: $\phi_\varepsilon^*g=e^{2\varepsilon}g$). \checkmark
\end{itemize}
\end{example}

\begin{intuition}[{why Killing fields matter in GR}]
Symmetries of spacetime give conserved quantities: a time-translation Killing field gives energy conservation, a rotation Killing field gives angular momentum (Section~\ref{sec:lc}). In the Schwarzschild solution $\partial_t$ and $\partial_\varphi$ are Killing.
\end{intuition}

% Section 13.3: Cartan's magic formula | Part: 2 | Depends on: sec:fields, sec:interior | Status: todo | Original lines: 2230-2254
%!TEX root = ../../../main.tex
\subsection{Cartan's magic formula}\label{sec:cartan}

\begin{theorem}[Cartan]\label{thm:cartan}
For every differential form $\omega$ and vector field $X$,
\begin{equation}\label{eq:cartan}
\Lie_X\omega=\dd\,\iota_X\omega+\iota_X\,\dd\omega .
\end{equation}
\end{theorem}
\begin{proof}
\Step{1}{both sides are derivations} $\Lie_X$ is a derivation of $\wedge$ (Leibniz rule for the flow derivative). The right-hand side is the graded commutator $[\dd,\iota_X]$ of two antiderivations (\eqref{eq:antider}, \eqref{eq:dleibniz}), hence a derivation of degree $0$.
\Step{2}{they commute with $\dd$} $\Lie_X\dd=\dd\Lie_X$ (pull-back commutes with $\dd$), and $\dd[\dd,\iota_X]=\dd\iota_X\dd=[\dd,\iota_X]\dd$ because $\dd\dd=0$.
\Step{3}{they agree on functions} $\dd\iota_Xf+\iota_X\dd f=0+X(f)=\Lie_Xf$.
\Step{4}{conclusion} Locally every form is a sum of terms $f\,\dd g_1\wedge\cdots\wedge\dd g_k$; two degree-$0$ derivations commuting with $\dd$ and agreeing on functions agree on $f$ and on $\dd g_i=\dd(g_i)$, hence on all forms.
\end{proof}

\begin{example}[Cartan's formula for a $1$-form]\label{ex:cartan1}
Take $\omega=x\,\dd y$ on $\R^2$ and $X=\partial_x$ (flow $\phi_t(x,y)=(x+t,y)$). Then $\iota_X\omega=x\,\dd y(\partial_x)=0$, so $\dd\iota_X\omega=0$; $\dd\omega=\dd x\wedge\dd y$ and $\iota_X\dd\omega=\dd y$. Hence $\dd\iota_X\omega+\iota_X\dd\omega=\dd y$.
Directly from the flow: $\phi_t\pb\omega=(x+t)\,\dd y$ and $\frac{\dd}{\dd t}\big|_0\phi_t\pb\omega=\dd y$ \checkmark.
In components, for $\omega=\omega_i\dd x^i$: $\Lie_X\omega=(X^j\partial_j\omega_i+\omega_j\partial_iX^j)\,\dd x^i$, and expanding $\dd(\omega_jX^j)+\iota_X(\partial_j\omega_i\,\dd x^j\wedge\dd x^i)$ gives the same expression.
\end{example}

\begin{remark}[Cartan's formula is the engine of the integral theorems]
Since $\dd\mu=0$ for a top-degree form, \eqref{eq:cartan} gives $\Lie_X\mu=\dd(\iota_X\mu)$: the rate of change of volume along a flow is the exterior derivative of the flux form $\iota_X\mu$. This is precisely the divergence theorem in disguise (Section~\ref{sec:gauss}); and for fluids it is the step that turns Euler's equation into the exact form $(\partial_t+\Lie_u)u^\flat=\dd(\cdots)$ behind Kelvin's circulation theorem (Section~\ref{sec:phys-fluids}).
\end{remark}

% Section 13: Operator relations: $[\partial_i,x^j]$ and $\{\iota_i,\dd x^j\wedge\}$} | Part: 2 | Depends on: see docs/REWRITE_PLAN.md | Status: todo | Original lines: 2255-2280
%!TEX root = ../../../main.tex
\subsection{Operator relations: $[\partial_i,x^j]$ and $\{\iota_i,\dd x^j\wedge\}$}\label{sec:operators}

Let $x^j$ act on forms by multiplication, $\partial_i$ by differentiating the coefficients, $\epsilon^j=\dd x^j\wedge$ by wedging, and $\iota_i=\iota_{\partial_i}$ by contracting.

\begin{proposition}\label{prop:operators}
For every form $\alpha$ and function $f$:
\begin{align}
\partial_i(x^jf)-x^j\partial_if&=\delta^j_if &&\text{i.e. }[\partial_i,x^j]=\delta^j_i,\label{eq:E12}\\
\iota_i(\dd x^j\wedge\alpha)+\dd x^j\wedge\iota_i\alpha&=\delta^j_i\alpha &&\text{i.e. }\{\iota_i,\dd x^j\wedge\}=\delta^j_i,\label{eq:E13}\\
\dd(x^j\alpha)-x^j\dd\alpha&=\dd x^j\wedge\alpha &&\text{i.e. }[\dd,x^j]=\dd x^j\wedge,\label{eq:E14}\\
\dd(\alpha_I\,\dd x^I)&=\partial_i\alpha_I\,\dd x^i\wedge\dd x^I &&\text{i.e. }\dd=\dd x^i\wedge\partial_i\text{ on coefficients,}\label{eq:E15}\\
\Lie_{\partial_i}(\alpha_I\,\dd x^I)&=(\dd\iota_i+\iota_i\dd)(\alpha_I\dd x^I)=\partial_i\alpha_I\,\dd x^I.\label{eq:E16}
\end{align}
\end{proposition}
\begin{proof}
\eqref{eq:E12}: Leibniz rule, $\partial_i(x^jf)=\delta^j_if+x^j\partial_if$. \eqref{eq:E13}: the antiderivation rule \eqref{eq:antider} with the $1$-form $\dd x^j$ gives $\iota_i(\dd x^j\wedge\alpha)=(\iota_i\dd x^j)\alpha-\dd x^j\wedge\iota_i\alpha=\delta^j_i\alpha-\dd x^j\wedge\iota_i\alpha$. \eqref{eq:E14}: Leibniz rule \eqref{eq:dleibniz} with a $0$-form, $\dd(x^j\alpha)=\dd x^j\wedge\alpha+x^j\dd\alpha$. \eqref{eq:E15} is the definition of $\dd$ in coordinates. For \eqref{eq:E16}: $\Lie_{\partial_i}\dd x^j=\dd(\iota_i\dd x^j)+\iota_i\dd\dd x^j=\dd(\delta^j_i)+0=0$, and $\Lie_{\partial_i}$ is a derivation, so only the coefficient is differentiated.
\end{proof}
\numcheck Check of \eqref{eq:E13} in $\R^2$ with $\alpha=x\,\dd y$ and $i=j=x$: $\dd x\wedge\alpha=x\,\dd x\wedge\dd y$ and $\iota_{\partial_x}(x\,\dd x\wedge\dd y)=x\,\dd y$; $\iota_{\partial_x}\alpha=x\,\dd y(\partial_x)=0$. The sum is $x\,\dd y=\alpha$ \checkmark.
Check of \eqref{eq:E16} with $\alpha=f\,\dd x\wedge\dd y$: $\dd\alpha=0$ (a $3$-form in the plane), so $\Lie_{\partial_x}\alpha=\dd(f\,\dd y)=f_x\,\dd x\wedge\dd y$ \checkmark.

\begin{remark}[Heisenberg and the fermions]
\eqref{eq:E12} and \eqref{eq:E13} are the same duality $\dd x^j(\partial_i)=\delta^j_i$ written for the even pair $(x^j,\partial_i)$ with a commutator and for the odd pair $(\dd x^j\wedge,\iota_i)$ with an anticommutator.
The first is the Heisenberg relation of quantum mechanics ($p=-i\hbar\partial_x$, $[x,p]=i\hbar$); the second is the canonical anticommutation relation of fermions, and the exterior algebra $\bigwedge(V^*)$ is the corresponding Fock space.
In both cases $\partial_x$ (resp.\ $\iota_{\partial_x}$) is the ``annihilation'' partner of $x$ (resp.\ $\dd x\wedge$).
\end{remark}

% Section 13 (keypoints): Keypoints | Part: 2 | Depends on: sec:fields, sec:interior | Status: todo | Original lines: 2281-2288
%!TEX root = ../../../main.tex
\begin{keypoints}[manifolds, forms and the Lie derivative]
\begin{itemize}
\item $x$ is a function, $\dd x=\dd(x)$ a covector, $\partial/\partial x$ the dual vector, with $\dd x(\partial_x)=1$. $\dd x$ depends only on the function $x$; $\partial/\partial x$ depends on the whole chart.
\item $\partial/\partial x=\iota_{\partial_x}\dd$ on functions; $\dd=\dd x^i\wedge\partial_i$ on forms; and $[\partial_i,x^j]=\{\iota_i,\dd x^j\wedge\}=\delta^j_i$.
\item The three derivatives of geometry: $\dd$ (forms; no extra data), $\Lie_X$ (needs a vector field), $\nabla$ (needs a connection); they are tied by $\Lie_X=\dd\iota_X+\iota_X\dd$ on forms.
\end{itemize}
\end{keypoints}

